\chapter{Data Structures --- Important Problems and Solutions for Recruitment Examinations}
\label{ch:data_structures_problems}

This chapter brings together essential practical programming questions, algorithm-tracing problems, and mathematical calculations from across the Data Structures syllabus. The problems are selected to reflect the testing patterns of Bangladesh government technical recruitment examinations, including Assistant Programmer, Assistant Director (ICT), Assistant Network Engineer, Assistant Maintenance Engineer, Assistant Database Administrator, and related IT cadres.

Each problem includes the problem statement, underlying algorithmic reasoning, complete worked solution, time and space complexity analysis, and common examination pitfalls.

\begin{center}
\vspace{0.5em}
{\color{primary}\rule{0.45\textwidth}{1.2pt}}\\[4pt]
{\Large\bfseries\color{primary} Section I: Arrays and Basic Algorithms}\\[2pt]
{\color{primary}\rule{0.45\textwidth}{1.2pt}}
\vspace{0.5em}
\end{center}

\section{Arrays and Basic Algorithms}

\subsection{Problem 1: Find the Smallest Element in an Array}

\textbf{Problem Statement:} Given an array of $n$ integers, find the smallest element without sorting the array.

\begin{lstlisting}[style=pseudocode, caption={Input}]
6
25 12 8 34 5 19
\end{lstlisting}

\begin{lstlisting}[style=pseudocode, caption={Expected Output}]
Smallest element: 5
\end{lstlisting}

\subsubsection*{Algorithmic Strategy}
Initialize the running minimum with the first element $A[0]$. Traverse the array from index $1$ to $n-1$, comparing each element against the running minimum. Whenever an element $A[i] < \text{minimum}$ is encountered, update $\text{minimum} = A[i]$.

\begin{lstlisting}[style=cppstyle, caption={C++ Solution for Finding Minimum Element.}]
#include <iostream>
#include <vector>
using namespace std;

int main() {
    int n;
    if (!(cin >> n) || n <= 0) {
        cout << "Array must not be empty\n";
        return 0;
    }

    vector<int> a(n);
    for (int i = 0; i < n; i++) {
        cin >> a[i];
    }

    int minimum = a[0];
    for (int i = 1; i < n; i++) {
        if (a[i] < minimum) {
            minimum = a[i];
        }
    }

    cout << "Smallest element: " << minimum << '\n';
    return 0;
}
\end{lstlisting}

\subsubsection*{Step-by-Step Execution Trace}
\begin{table}[htbp]
\centering
\small
\begin{tabular}{c c c}
\toprule
\textbf{Step} & \textbf{Element Examined ($A[i]$)} & \textbf{Current Minimum} \\
\midrule
$i = 0$ & 25 (Initial) & 25 \\
$i = 1$ & 12 ($12 < 25$) & 12 \\
$i = 2$ & 8 ($8 < 12$) & 8 \\
$i = 3$ & 34 ($34 \geq 8$) & 8 \\
$i = 4$ & 5 ($5 < 8$) & 5 \\
$i = 5$ & 19 ($19 \geq 5$) & 5 \\
\bottomrule
\end{tabular}
\caption{Linear scan trace for finding the minimum element.}
\end{table}

\textbf{Complexity Analysis:}
\begin{itemize}
    \item \textbf{Time Complexity:} $O(n)$, since every element is examined exactly once.
    \item \textbf{Auxiliary Space:} $O(1)$, requiring only a single scalar variable.
\end{itemize}

\begin{quote}
\textbf{Exam Note:} Sorting the array first takes $O(n \log n)$ time and is unnecessary. Linear scanning in $O(n)$ is asymptotically optimal because any adversary could place the minimum in the unexamined position.
\end{quote}

\subsection{Problem 2: Find the Second-Largest Distinct Element}

\textbf{Problem Statement:} Find the second-largest \textbf{distinct} value in an integer array without sorting it.

\begin{lstlisting}[style=pseudocode, caption={Input}]
7
10 25 8 25 17 30 30
\end{lstlisting}

\begin{lstlisting}[style=pseudocode, caption={Expected Output}]
Second largest: 25
\end{lstlisting}

\subsubsection*{Algorithmic Strategy}
Maintain two variables: \texttt{largest} and \texttt{secondLargest}, initialized to $-\infty$.
\begin{itemize}
    \item If $A[i] > \text{largest}$: update $\text{secondLargest} = \text{largest}$, then $\text{largest} = A[i]$.
    \item Else if $A[i] < \text{largest}$ and $A[i] > \text{secondLargest}$: update $\text{secondLargest} = A[i]$.
    \item If $A[i] == \text{largest}$: ignore duplicates of the largest value.
\end{itemize}

\begin{lstlisting}[style=cppstyle, caption={C++ Solution for Finding Second-Largest Distinct Element.}]
#include <iostream>
#include <vector>
#include <climits>
using namespace std;

int main() {
    int n;
    if (!(cin >> n) || n < 2) {
        cout << "Array must contain at least two elements\n";
        return 0;
    }

    vector<int> a(n);
    for (int i = 0; i < n; i++) {
        cin >> a[i];
    }

    long long largest = LLONG_MIN;
    long long secondLargest = LLONG_MIN;
    bool hasLargest = false;
    bool hasSecond = false;

    for (int x : a) {
        if (!hasLargest || x > largest) {
            if (hasLargest) {
                secondLargest = largest;
                hasSecond = true;
            }
            largest = x;
            hasLargest = true;
        } else if (x < largest && (!hasSecond || x > secondLargest)) {
            secondLargest = x;
            hasSecond = true;
        }
    }

    if (hasSecond) {
        cout << "Second largest: " << secondLargest << '\n';
    } else {
        cout << "No second-largest distinct value\n";
    }
    return 0;
}
\end{lstlisting}

\textbf{Complexity Analysis:}
\begin{itemize}
    \item \textbf{Time Complexity:} $O(n)$ in a single pass.
    \item \textbf{Auxiliary Space:} $O(1)$.
\end{itemize}

\subsection{Problem 3: Binary Search on a Sorted Array}

\textbf{Problem Statement:} Search for target value $T = 23$ in the sorted array:
\[
A = [3, 8, 12, 17, 23, 31, 42]
\]

\subsubsection*{Algorithmic Strategy}
Maintain search boundaries $[\text{left}, \text{right}]$. Compute midpoint:
\[
\text{mid} = \text{left} + \left\lfloor \frac{\text{right} - \text{left}}{2} \right\rfloor
\]
\begin{enumerate}
    \item Initially: $\text{left} = 0, \text{right} = 6 \implies \text{mid} = 3$, $A[3] = 17$. Since $23 > 17$, eliminate left half: $\text{left} = 4$.
    \item New interval: $[4, 6] \implies \text{mid} = 4 + 1 = 5$, $A[5] = 31$. Since $23 < 31$, eliminate right half: $\text{right} = 4$.
    \item New interval: $[4, 4] \implies \text{mid} = 4$, $A[4] = 23$. Target found at index $4$!
\end{enumerate}

\begin{lstlisting}[style=cppstyle, caption={C++ Iterative Binary Search.}]
#include <iostream>
#include <vector>
using namespace std;

int binarySearch(const vector<int>& a, int target) {
    int left = 0;
    int right = static_cast<int>(a.size()) - 1;

    while (left <= right) {
        int mid = left + (right - left) / 2; // Prevents integer overflow
        if (a[mid] == target) return mid;
        if (a[mid] < target)  left = mid + 1;
        else                  right = mid - 1;
    }
    return -1; // Target not found
}

int main() {
    vector<int> a = {3, 8, 12, 17, 23, 31, 42};
    int index = binarySearch(a, 23);
    cout << "Target 23 found at index: " << index << '\n';
    return 0;
}
\end{lstlisting}

\textbf{Complexity Analysis:}
\begin{itemize}
    \item \textbf{Time Complexity:} $O(\log n)$ worst and average case; $O(1)$ best case.
    \item \textbf{Auxiliary Space:} $O(1)$ for iterative; $O(\log n)$ for recursive call stack.
\end{itemize}

\begin{quote}
\textbf{Exam Trap:} Binary search \textbf{strictly requires} a sorted collection. It cannot be applied directly to an unsorted array without prior sorting.
\end{quote}

\newpage
\begin{center}
\vspace{0.5em}
{\color{primary}\rule{0.45\textwidth}{1.2pt}}\\[4pt]
{\Large\bfseries\color{primary} Section II: Stack and Queue Problems}\\[2pt]
{\color{primary}\rule{0.45\textwidth}{1.2pt}}
\vspace{0.5em}
\end{center}

\section{Stack and Queue Problems}

\subsection{Problem 4: Balanced Parentheses Checking}

\textbf{Problem Statement:} Determine whether the following bracket expressions are balanced:
\begin{itemize}
    \item Expression 1: \texttt{\{[a + (b * c)]\}} $\implies$ \textbf{Balanced}
    \item Expression 2: \texttt{\{[a + (b * c]\}} $\implies$ \textbf{Not balanced}
\end{itemize}

\subsubsection*{Algorithmic Strategy}
Use a stack to record opening brackets:
\begin{enumerate}
    \item When an opening delimiter (\texttt{(}, \texttt{\{}, \texttt{[}) is read, push it onto the stack.
    \item When a closing delimiter (\texttt{)}, \texttt{\}}, \texttt{]}) is read:
    \begin{itemize}
        \item If stack is empty, return false (unmatched closing bracket).
        \item Pop top element; if it does not match the closing bracket's type, return false.
    \end{itemize}
    \item At the end of the string, the expression is balanced if and only if the stack is completely empty.
\end{enumerate}

\begin{lstlisting}[style=cppstyle, caption={C++ Solution for Balanced Parentheses.}]
#include <iostream>
#include <stack>
#include <string>
using namespace std;

bool isMatching(char opening, char closing) {
    return (opening == '(' && closing == ')') ||
           (opening == '[' && closing == ']') ||
           (opening == '{' && closing == '}');
}

bool isBalanced(const string& expression) {
    stack<char> st;
    for (char ch : expression) {
        if (ch == '(' || ch == '[' || ch == '{') {
            st.push(ch);
        } else if (ch == ')' || ch == ']' || ch == '}') {
            if (st.empty() || !isMatching(st.top(), ch)) {
                return false;
            }
            st.pop();
        }
    }
    return st.empty();
}

int main() {
    string exp1 = "{[a + (b * c)]}";
    string exp2 = "{[a + (b * c]}";

    cout << exp1 << " -> " << (isBalanced(exp1) ? "Balanced" : "Not balanced") << '\n';
    cout << exp2 << " -> " << (isBalanced(exp2) ? "Balanced" : "Not balanced") << '\n';
    return 0;
}
\end{lstlisting}

\textbf{Complexity Analysis:}
\begin{itemize}
    \item \textbf{Time Complexity:} $O(n)$ where $n$ is string length.
    \item \textbf{Auxiliary Space:} $O(n)$ in the worst case (e.g., all opening brackets).
\end{itemize}

\subsection{Problem 5: Postfix Expression Evaluation}

\textbf{Problem Statement:} Evaluate the postfix (Reverse Polish Notation) expression:
\[
5 \quad 3 \quad 2 \quad * \quad +
\]

\subsubsection*{Step-by-Step Evaluation Trace}
\begin{enumerate}
    \item Push operand $5$: Stack $= [5]$.
    \item Push operand $3$: Stack $= [5, 3]$.
    \item Push operand $2$: Stack $= [5, 3, 2]$.
    \item Operator $*$ encountered:
    \begin{itemize}
        \item Pop $\text{op}_2 = 2$, Pop $\text{op}_1 = 3$.
        \item Compute: $\text{result} = 3 \times 2 = 6$.
        \item Push $6$: Stack $= [5, 6]$.
    \end{itemize}
    \item Operator $+$ encountered:
    \begin{itemize}
        \item Pop $\text{op}_2 = 6$, Pop $\text{op}_1 = 5$.
        \item Compute: $\text{result} = 5 + 6 = 11$.
        \item Push $11$: Stack $= [11]$.
    \end{itemize}
\end{enumerate}
Final result:
\[
\boxed{11}
\]

\begin{quote}
\textbf{Crucial Rule for Non-Commutative Operators:} In division and subtraction, order is vital:
\[
\text{op}_2 = \text{pop}(), \quad \text{op}_1 = \text{pop}() \implies \text{compute } (\text{op}_1 - \text{op}_2) \text{ or } (\text{op}_1 / \text{op}_2).
\]
\end{quote}

\subsection{Problem 6: Implement a Queue Using Two Stacks}

\textbf{Problem Statement:} Implement FIFO queue semantics using two LIFO stacks. Trace:
\[
\text{enqueue}(10), \; \text{enqueue}(20), \; \text{enqueue}(30), \; \text{dequeue}(), \; \text{enqueue}(40), \; \text{dequeue}()
\]

\begin{lstlisting}[style=cppstyle, caption={C++ Implementation of Queue using Two Stacks.}]
#include <iostream>
#include <stack>
using namespace std;

class QueueUsingStacks {
private:
    stack<int> inStack;
    stack<int> outStack;

    void transferIfNeeded() {
        if (outStack.empty()) {
            while (!inStack.empty()) {
                outStack.push(inStack.top());
                inStack.pop();
            }
        }
    }

public:
    void enqueue(int value) {
        inStack.push(value);
    }

    bool dequeue(int& value) {
        transferIfNeeded();
        if (outStack.empty()) return false;
        value = outStack.top();
        outStack.pop();
        return true;
    }
};

int main() {
    QueueUsingStacks q;
    int val;

    q.enqueue(10);
    q.enqueue(20);
    q.enqueue(30);

    q.dequeue(val);
    cout << "Dequeued: " << val << '\n'; // 10

    q.enqueue(40);

    q.dequeue(val);
    cout << "Dequeued: " << val << '\n'; // 20

    return 0;
}
\end{lstlisting}

\textbf{Complexity Analysis:}
\begin{itemize}
    \item \textbf{Enqueue:} $O(1)$ worst case.
    \item \textbf{Dequeue:} $O(1)$ amortized; $O(n)$ worst case (when \texttt{outStack} is empty).
    \item \textbf{Space:} $O(n)$ total storage across both stacks.
\end{itemize}

\newpage
\begin{center}
\vspace{0.5em}
{\color{primary}\rule{0.45\textwidth}{1.2pt}}\\[4pt]
{\Large\bfseries\color{primary} Section III: Linked-List Problems}\\[2pt]
{\color{primary}\rule{0.45\textwidth}{1.2pt}}
\vspace{0.5em}
\end{center}

\section{Linked-List Problems}

\subsection{Problem 7: Reverse a Singly Linked List}

\textbf{Problem Statement:} Reverse the list $10 \to 20 \to 30 \to 40 \to \text{NULL}$ in place.

\subsubsection*{Algorithmic Strategy}
Maintain three pointers: \texttt{previous = nullptr}, \texttt{current = head}, and \texttt{nextNode = nullptr}.
\begin{enumerate}
    \item Save next node: $\texttt{nextNode} = \texttt{current->next}$.
    \item Reverse current pointer: $\texttt{current->next} = \texttt{previous}$.
    \item Advance previous: $\texttt{previous} = \texttt{current}$.
    \item Advance current: $\texttt{current} = \texttt{nextNode}$.
\end{enumerate}

\begin{lstlisting}[style=cppstyle, caption={C++ Solution for In-Place Linked List Reversal.}]
#include <iostream>
using namespace std;

struct Node {
    int data;
    Node* next;
    explicit Node(int val) : data(val), next(nullptr) {}
};

Node* reverseList(Node* head) {
    Node* previous = nullptr;
    Node* current = head;

    while (current != nullptr) {
        Node* nextNode = current->next;
        current->next = previous;
        previous = current;
        current = nextNode;
    }
    return previous; // New head
}

void display(Node* head) {
    while (head) {
        cout << head->data << (head->next ? " -> " : " -> NULL\n");
        head = head->next;
    }
}

int main() {
    Node* head = new Node(10);
    head->next = new Node(20);
    head->next->next = new Node(30);
    head->next->next->next = new Node(40);

    cout << "Original list: "; display(head);
    head = reverseList(head);
    cout << "Reversed list: "; display(head);

    return 0;
}
\end{lstlisting}

\textbf{Complexity Analysis:} Time $O(n)$; Auxiliary space $O(1)$.

\subsection{Problem 8: Detect a Cycle in a Linked List (Floyd's Algorithm)}

\textbf{Problem Statement:} Detect if a singly linked list contains a cycle using $O(1)$ auxiliary memory.

\subsubsection*{Floyd's Tortoise and Hare Algorithm}
Use two pointers:
\begin{itemize}
    \item \texttt{slow}: advances $1$ node per step.
    \item \texttt{fast}: advances $2$ nodes per step.
\end{itemize}
If a cycle exists, the relative distance between them decreases by $1$ in each step inside the loop, guaranteeing that \texttt{fast} will catch up with \texttt{slow} in at most $L$ steps (where $L$ is loop length). If no cycle exists, \texttt{fast} reaches \texttt{NULL}.

\begin{lstlisting}[style=cppstyle, caption={C++ Cycle Detection using Floyd's Algorithm.}]
bool hasCycle(Node* head) {
    Node* slow = head;
    Node* fast = head;

    while (fast != nullptr && fast->next != nullptr) {
        slow = slow->next;
        fast = fast->next->next;
        if (slow == fast) {
            return true; // Cycle detected
        }
    }
    return false; // Reached end: no cycle
}
\end{lstlisting}

\textbf{Complexity Analysis:} Time $O(n)$; Auxiliary space $O(1)$.

\newpage
\begin{center}
\vspace{0.5em}
{\color{primary}\rule{0.45\textwidth}{1.2pt}}\\[4pt]
{\Large\bfseries\color{primary} Section IV: Heap and Priority-Queue Problems}\\[2pt]
{\color{primary}\rule{0.45\textwidth}{1.2pt}}
\vspace{0.5em}
\end{center}

\section{Heap and Priority-Queue Problems}

\subsection{Problem 9: Build a Max-Heap from an Unordered Array}

\textbf{Problem Statement:} Construct a max-heap from array $A = [4, 10, 3, 5, 1]$.

\subsubsection*{Worked Solution}
For array of size $n = 5$, the last non-leaf node index is:
\[
\left\lfloor \frac{n}{2} \right\rfloor - 1 = \left\lfloor \frac{5}{2} \right\rfloor - 1 = 1
\]
Heapify in reverse order from index $1$ down to $0$:
\begin{itemize}
    \item \textbf{At index 1 ($A[1] = 10$):} Children are $A[3] = 5, A[4] = 1$. Since $10 > 5$ and $10 > 1$, max-heap property is already satisfied.
    \item \textbf{At index 0 ($A[0] = 4$):} Children are $A[1] = 10, A[2] = 3$. Largest child is $10$. Swap $A[0]$ and $A[1]$:
    \[
    A = [10, 4, 3, 5, 1]
    \]
    Sift down $4$ at index $1$: Children are $A[3] = 5, A[4] = 1$. Largest is $5$. Swap $A[1]$ and $A[3]$:
    \[
    \boxed{A = [10, 5, 3, 4, 1]}
    \]
\end{itemize}

\begin{quote}
\textbf{Key Complexity Theorem:} Building a heap bottom-up takes $\boxed{O(n)}$ time, not $O(n \log n)$, because the sum of heights across all nodes converges to $O(n)$:
\[
\sum_{h=0}^{\lfloor \log n \rfloor} \frac{n}{2^{h+1}} O(h) = O(n).
\]
\end{quote}

\subsection{Problem 10: Insert and Delete Maximum from a Max-Heap}

\textbf{Problem Statement:} Given max-heap $A = [10, 5, 3, 4, 1]$:
\begin{enumerate}
    \item Insert element $15$.
    \item Delete the maximum element.
\end{enumerate}

\subsubsection*{1. Insertion Trace}
Append $15$ at index $5$: $A = [10, 5, 3, 4, 1, 15]$.
\begin{itemize}
    \item Parent of index $5$ is $\lfloor (5-1)/2 \rfloor = 2$, holding value $3$. Since $15 > 3$, swap:
    \[
    A = [10, 5, 15, 4, 1, 3]
    \]
    \item Parent of index $2$ is $\lfloor (2-1)/2 \rfloor = 0$, holding value $10$. Since $15 > 10$, swap:
    \[
    \boxed{A = [15, 5, 10, 4, 1, 3]}
    \]
\end{itemize}

\subsubsection*{2. Deletion Trace}
Extract root $15$. Move last element ($3$) to root and shrink heap size:
\[
A = [3, 5, 10, 4, 1]
\]
Sift down root $3$: Children are $5$ (idx 1) and $10$ (idx 2). Larger child is $10$. Swap $3$ and $10$:
\[
\boxed{A = [10, 5, 3, 4, 1]}
\]

\subsection{Problem 11: Priority Queue Extraction Order}

\textbf{Problem Statement:} A max-priority queue contains elements $\{12, 5, 20, 8, 15\}$. In what sequence are they extracted?
\begin{itemize}
    \item \textbf{Max-Priority Queue:} Extracts in descending order of priority:
    \[
    20 \longrightarrow 15 \longrightarrow 12 \longrightarrow 8 \longrightarrow 5
    \]
    \item \textbf{Min-Priority Queue:} Extracts in ascending order:
    \[
    5 \longrightarrow 8 \longrightarrow 12 \longrightarrow 15 \longrightarrow 20
    \]
\end{itemize}

\newpage
\begin{center}
\vspace{0.5em}
{\color{primary}\rule{0.45\textwidth}{1.2pt}}\\[4pt]
{\Large\bfseries\color{primary} Section V: Circular Queue Problems}\\[2pt]
{\color{primary}\rule{0.45\textwidth}{1.2pt}}
\vspace{0.5em}
\end{center}

\section{Circular Queue Problems}

\subsection{Problem 12: Trace Circular Queue Operations}

\textbf{Problem Statement:} A circular queue has buffer capacity $N = 5$. Initially, $\texttt{front} = \texttt{rear} = -1$. Trace the following operations and determine if the final queue is full:
\[
\text{enq}(10), \text{enq}(20), \text{enq}(30), \text{enq}(40), \text{deq}(), \text{deq}(), \text{enq}(50), \text{enq}(60), \text{enq}(70)
\]

\subsubsection*{Step-by-Step Execution Trace}
Index advance formula: $\text{next} = (\text{index} + 1) \bmod N$. Full condition: $(\text{rear} + 1) \bmod N == \text{front}$.

\begin{table}[htbp]
\centering
\small
\begin{tabularx}{\textwidth}{l c c X}
\toprule
\textbf{Operation} & \textbf{Front} & \textbf{Rear} & \textbf{Logical Queue Content} \\
\midrule
Initial State & -1 & -1 & Queue is Empty \\
$\text{enqueue}(10)$ & 0 & 0 & $[10]$ \\
$\text{enqueue}(20)$ & 0 & 1 & $[10, 20]$ \\
$\text{enqueue}(30)$ & 0 & 2 & $[10, 20, 30]$ \\
$\text{enqueue}(40)$ & 0 & 3 & $[10, 20, 30, 40]$ \\
$\text{dequeue}()$   & 1 & 3 & $[20, 30, 40]$ ($10$ removed) \\
$\text{dequeue}()$   & 2 & 3 & $[30, 40]$ ($20$ removed) \\
$\text{enqueue}(50)$ & 2 & 4 & $[30, 40, 50]$ \\
$\text{enqueue}(60)$ & 2 & 0 & $[30, 40, 50, 60]$ (Wrapped to idx 0) \\
$\text{enqueue}(70)$ & 2 & 1 & $[30, 40, 50, 60, 70]$ (Wrapped to idx 1) \\
\bottomrule
\end{tabularx}
\caption{Trace of circular queue operations.}
\end{table}

\textbf{Final Fullness Evaluation:}
\[
(\text{rear} + 1) \bmod 5 = (1 + 1) \bmod 5 = 2 == \text{front}
\]
The queue contains 5 elements and is \textbf{strictly full}. Any further insertion triggers queue overflow.

\newpage
\begin{center}
\vspace{0.5em}
{\color{primary}\rule{0.45\textwidth}{1.2pt}}\\[4pt]
{\Large\bfseries\color{primary} Section VI: Tree Problems}\\[2pt]
{\color{primary}\rule{0.45\textwidth}{1.2pt}}
\vspace{0.5em}
\end{center}

\section{Tree Problems}

\subsection{Problem 13: Tree Traversal Orders}

\textbf{Problem Statement:} Consider the Binary Search Tree:
\begin{itemize}
    \item Root: $8$. Left child: $3$. Right child: $10$.
    \item Children of $3$: left $1$, right $6$. Children of $6$: left $4$, right $7$.
    \item Children of $10$: right $14$. Child of $14$: left $13$.
\end{itemize}
Determine the four standard traversals.

\subsubsection*{Solutions}
\begin{itemize}
    \item \textbf{Preorder (Root, Left, Right):}
    \[
    8, \; 3, \; 1, \; 6, \; 4, \; 7, \; 10, \; 14, \; 13
    \]
    \item \textbf{Inorder (Left, Root, Right):}
    \[
    \boxed{1, \; 3, \; 4, \; 6, \; 7, \; 8, \; 10, \; 13, \; 14} \quad (\text{strictly sorted in ascending order})
    \]
    \item \textbf{Postorder (Left, Right, Root):}
    \[
    1, \; 4, \; 7, \; 6, \; 3, \; 13, \; 14, \; 10, \; 8
    \]
    \item \textbf{Level-Order (Breadth-First):}
    \[
    8, \; 3, \; 10, \; 1, \; 6, \; 14, \; 4, \; 7, \; 13
    \]
\end{itemize}

\subsection{Problem 14: Leaves, Nodes, and Edges in a Full Binary Tree}

\textbf{Problem Statement:} A full binary tree has $I = 20$ internal nodes. Calculate leaves $L$, total nodes $N$, and total edges $E$.

\subsubsection*{Worked Calculation}
For any full binary tree (where every non-leaf node has exactly 2 children):
\[
L = I + 1 = 20 + 1 = \boxed{21 \text{ leaves}}
\]
\[
N = I + L = 20 + 21 = \boxed{41 \text{ total nodes}}
\]
\[
E = N - 1 = 41 - 1 = \boxed{40 \text{ edges}} \quad (\text{also verified by } 2I = 2 \times 20 = 40)
\]

\subsection{Problem 15: Minimum Height of a Binary Tree}

\textbf{Problem Statement:} Find the minimum possible height (in edges) of a binary tree containing $N = 100$ nodes.

\subsubsection*{Worked Calculation}
A binary tree of height $h$ edges has maximum node capacity:
\[
N_{\max} = 2^{h+1} - 1
\]
For $N = 100$:
\[
100 \leq 2^{h+1} - 1 \implies 101 \leq 2^{h+1}
\]
Since $2^6 = 64$ and $2^7 = 128$:
\[
h + 1 \geq 7 \implies \boxed{h_{\min} = 6 \text{ edges}}
\]

\subsection{Problem 16: Maximum Data Keys in a B+ Tree}

\textbf{Problem Statement:} A B+ tree has order $m = 4$ and height $h = 3$ edges. Each internal node can have at most $4$ children, and each leaf can hold at most $3$ data keys. Calculate maximum data-key capacity.

\subsubsection*{Worked Calculation}
The maximum number of leaf nodes at height $h = 3$ is:
\[
L_{\max} = m^h = 4^3 = 64 \text{ leaf nodes}
\]
Each leaf holds up to $m - 1 = 3$ data records. Total maximum data records:
\[
\boxed{N_{\max} = 64 \times 3 = 192 \text{ records}}
\]

\newpage
\begin{center}
\vspace{0.5em}
{\color{primary}\rule{0.45\textwidth}{1.2pt}}\\[4pt]
{\Large\bfseries\color{primary} Section VII: Graph Problems}\\[2pt]
{\color{primary}\rule{0.45\textwidth}{1.2pt}}
\vspace{0.5em}
\end{center}

\section{Graph Problems}

\subsection{Problem 17: Edge Calculation using Handshaking Lemma}

\textbf{Problem Statement:} An undirected graph has five vertices with degrees: $2, 3, 2, 3, 2$. Find the number of edges.

\subsubsection*{Worked Calculation}
By the Handshaking Lemma:
\[
\sum_{v \in V} \deg(v) = 2E \implies 2 + 3 + 2 + 3 + 2 = 12
\]
\[
2E = 12 \implies \boxed{E = 6 \text{ edges}}
\]

\subsection{Problem 18: Maximum Edges in Simple Graphs}

\textbf{Problem Statement:} Calculate maximum edges for $V = 10$ vertices:
\begin{enumerate}
    \item Simple undirected graph:
    \[
    E_{\max} = \frac{V(V-1)}{2} = \frac{10 \times 9}{2} = \boxed{45 \text{ edges}}
    \]
    \item Simple directed graph without self-loops:
    \[
    E_{\max} = V(V-1) = 10 \times 9 = \boxed{90 \text{ arcs}}
    \]
\end{enumerate}

\subsection{Problem 19: Complete Bipartite Graph Edges}

\textbf{Problem Statement:} Calculate total edges in complete bipartite graph $K_{4,6}$.
\[
E = a \times b = 4 \times 6 = \boxed{24 \text{ edges}}
\]

\subsection{Problem 20: BFS Traversal and Shortest Path}

\textbf{Problem Statement:} An undirected graph has adjacency lists:
\begin{lstlisting}[style=pseudocode]
A: B, C
B: A, D, E
C: A, F
D: B
E: B, F
F: C, E
\end{lstlisting}
Perform BFS starting from $A$ and find the shortest path from $A$ to $F$.

\subsubsection*{Execution Trace}
\begin{itemize}
    \item Step 1: Enqueue $A$, mark visited: $\text{Queue} = [A]$.
    \item Step 2: Dequeue $A$; enqueue unvisited neighbors $B, C$: $\text{Queue} = [B, C]$.
    \item Step 3: Dequeue $B$; enqueue unvisited neighbors $D, E$: $\text{Queue} = [C, D, E]$.
    \item Step 4: Dequeue $C$; enqueue unvisited neighbor $F$: $\text{Queue} = [D, E, F]$.
    \item Step 5: Dequeue $D, E, F$.
\end{itemize}
\textbf{BFS Visit Order:} $\boxed{A, B, C, D, E, F}$.\\
\textbf{Shortest Path to $F$:} $\boxed{A \to C \to F}$ (length: 2 edges).

\newpage
\begin{center}
\vspace{0.5em}
{\color{primary}\rule{0.45\textwidth}{1.2pt}}\\[4pt]
{\Large\bfseries\color{primary} Section VIII: Hash Table Problems}\\[2pt]
{\color{primary}\rule{0.45\textwidth}{1.2pt}}
\vspace{0.5em}
\end{center}

\section{Hash Table Problems}

\subsection{Problem 21: Division Method Hash Index Calculation}

\textbf{Problem Statement:} Table size $m = 11$, $h(k) = k \bmod 11$. Calculate slot indices for keys: $32, 44, 55, 67, 78$.

\begin{table}[htbp]
\centering
\small
\begin{tabular}{c c c l}
\toprule
\textbf{Key ($k$)} & \textbf{Calculation ($k \bmod 11$)} & \textbf{Slot Index} & \textbf{Collision Status} \\
\midrule
32 & $32 \bmod 11$ & 10 & No collision \\
44 & $44 \bmod 11$ & 0  & First key at slot 0 \\
55 & $55 \bmod 11$ & 0  & \textbf{Collides with 44 at slot 0} \\
67 & $67 \bmod 11$ & 1  & First key at slot 1 \\
78 & $78 \bmod 11$ & 1  & \textbf{Collides with 67 at slot 1} \\
\bottomrule
\end{tabular}
\caption{Hash indices and collisions under division hashing.}
\end{table}

\subsection{Problem 22: Linear Probing Insertion}

\textbf{Problem Statement:} Table size $m = 10$, $h(k) = k \bmod 10$. Insert keys in order: $23, 43, 13, 27, 33$ using linear probing.

\subsubsection*{Step-by-Step Insertion}
\begin{itemize}
    \item Key 23: $23 \bmod 10 = 3$. Slot 3 empty $\implies$ store \textbf{23 at index 3}.
    \item Key 43: $43 \bmod 10 = 3$. Slot 3 occupied $\implies$ probe slot 4 $\implies$ store \textbf{43 at index 4}.
    \item Key 13: $13 \bmod 10 = 3$. Slots 3, 4 occupied $\implies$ probe slot 5 $\implies$ store \textbf{13 at index 5}.
    \item Key 27: $27 \bmod 10 = 7$. Slot 7 empty $\implies$ store \textbf{27 at index 7}.
    \item Key 33: $33 \bmod 10 = 3$. Slots 3, 4, 5 occupied $\implies$ probe slot 6 $\implies$ store \textbf{33 at index 6}.
\end{itemize}

\begin{table}[htbp]
\centering
\small
\begin{tabular}{c c c c c c c c c c c}
\toprule
\textbf{Index} & 0 & 1 & 2 & 3 & 4 & 5 & 6 & 7 & 8 & 9 \\
\midrule
\textbf{Content} & --- & --- & --- & \textbf{23} & \textbf{43} & \textbf{13} & \textbf{33} & \textbf{27} & --- & --- \\
\bottomrule
\end{tabular}
\caption{Final state of hash table after linear probing.}
\end{table}

\subsection{Problem 23: Load Factor Calculation}

\textbf{Problem Statement:} Table has $m = 50$ slots and contains $n = 35$ entries.
\[
\alpha = \frac{n}{m} = \frac{35}{50} = 0.70 \implies \boxed{70\% \text{ occupancy}}
\]

\subsection{Problem 24: Rehashing Threshold Calculation}

\textbf{Problem Statement:} Table capacity $m = 200$. Rehashing triggers when $\alpha > 0.75$. What is the maximum number of items stored before rehashing occurs?
\[
\frac{n}{200} \leq 0.75 \implies n \leq 200 \times 0.75 \implies \boxed{n_{\max} = 150 \text{ entries}}
\]
Inserting the $151^{\text{st}}$ item raises $\alpha = 151/200 = 0.755 > 0.75$, triggering table doubling and rehashing.

\newpage
\begin{center}
\vspace{0.5em}
{\color{primary}\rule{0.45\textwidth}{1.2pt}}\\[4pt]
{\Large\bfseries\color{primary} Section IX: Representation and Architecture Selection}\\[2pt]
{\color{primary}\rule{0.45\textwidth}{1.2pt}}
\vspace{0.5em}
\end{center}

\section{Graph Representation and Architecture Selection}

\subsection{Problem 25: Memory Comparison of Matrix vs List}

\textbf{Problem Statement:} An undirected graph has $V = 6$ vertices and $E = 7$ edges. Compare raw structural cell counts:
\begin{itemize}
    \item \textbf{Adjacency Matrix:} $V^2 = 6^2 = \boxed{36 \text{ entries}}$.
    \item \textbf{Adjacency List:} $V$ head pointers $+ 2E$ edge nodes $= 6 + 14 = \boxed{20 \text{ structural entries}}$.
\end{itemize}

\subsection{Problem 26: Master Data Structure Selection Guide}

\begin{table}[htbp]
\centering
\small
\begin{tabularx}{\textwidth}{p{7.5cm} l X}
\toprule
\textbf{Engineering Requirement} & \textbf{Optimal Structure} & \textbf{Primary Theoretical Justification} \\
\midrule
Undo operation in text editors & Stack & Strict LIFO retrieval \\
Print job scheduling / CPU task queue & Queue & Fair FIFO order processing \\
Fixed-size streaming packet buffer & Circular Queue & $O(1)$ slot recycling without shifting \\
Extract highest-urgency emergency task & Max-Priority Queue & Highest-key service first \\
Fast insertion/deletion without contiguous memory & Linked List & Dynamic pointer rearrangement \\
Retrieve maximum key repeatedly & Max-Heap & $O(1)$ root peek, $O(\log n)$ extract \\
Organizational hierarchy / File directory & Tree & Parent-child non-cyclic branching \\
Database indexing with range queries & B+ Tree & Massive disk fanout; linked leaves \\
Dictionary with expected $O(1)$ search & Hash Table & Direct key-to-slot mapping \\
Graph traversal on sparse networks & Adjacency List & $O(V+E)$ space; optimal neighbor loop \\
Repeated edge-existence test in dense graph & Adjacency Matrix & $O(1)$ direct array index lookup \\
Shortest path in unweighted graph & BFS (Queue) & Explores vertices by edge-hop levels \\
Cycle detection / Topological sorting & DFS (Stack) & Explores deepest paths before backtrack \\
Prefix matching / Auto-complete & Trie & Reuses shared string prefixes \\
\bottomrule
\end{tabularx}
\caption{Master data structure selection chart for technical examinations.}
\end{table}

\newpage
\begin{center}
\vspace{0.5em}
{\color{primary}\rule{0.45\textwidth}{1.2pt}}\\[4pt]
{\Large\bfseries\color{primary} Section X: Rapid-Fire Revision Questions}\\[2pt]
{\color{primary}\rule{0.45\textwidth}{1.2pt}}
\vspace{0.5em}
\end{center}

\section{Rapid-Fire Revision Questions}

\begin{longtable}{r p{9.8cm} p{6.8cm}}
\toprule
\textbf{\#} & \textbf{Recruitment Examination Question} & \textbf{Concise Technical Answer} \\
\midrule
\endfirsthead
\toprule
\textbf{\#} & \textbf{Recruitment Examination Question} & \textbf{Concise Technical Answer} \\
\midrule
\endhead
\bottomrule
\endfoot
\bottomrule
\endlastfoot
1  & Time complexity of finding minimum in unsorted array? & $O(n)$ \\
2  & Time complexity of iterative binary search? & $O(\log n)$ \\
3  & Which principle does a stack follow? & LIFO (Last In, First Out) \\
4  & Which principle does a queue follow? & FIFO (First In, First Out) \\
5  & Purpose of two stacks in queue implementation? & Order reversal to achieve FIFO semantics \\
6  & Time complexity of iterative singly linked list reversal? & $O(n)$ time, $O(1)$ auxiliary space \\
7  & Algorithm detecting linked list cycle with $O(1)$ space? & Floyd's Cycle-Finding (Tortoise and Hare) \\
8  & Time complexity of bottom-up heap construction? & $O(n)$ \\
9  & Location of maximum element in a max-heap? & At the root ($A[0]$) \\
10 & Worst-case insertion complexity in binary heap? & $O(\log n)$ \\
11 & What does a circular queue reuse? & Vacated front array slots via modulo \\
12 & What determines removal order in priority queue? & Priority key value \\
13 & Number of edges in a tree with $n$ vertices? & $n - 1$ \\
14 & Maximum nodes in binary tree of height $h$ edges? & $2^{h+1} - 1$ \\
15 & Relationship between leaves and internal nodes in full binary tree? & $L = I + 1$ \\
16 & Height of skewed binary tree with $n$ nodes? & $n - 1$ edges \\
17 & Worst-case search complexity of unbalanced BST? & $O(n)$ \\
18 & Which traversal outputs keys of BST in sorted order? & Inorder traversal \\
19 & AVL tree balance factor restriction? & $|BF| \leq 1$ (i.e., $-1, 0, +1$) \\
20 & Maximum children in order-$m$ B-tree? & $m$ \\
21 & Where are actual data records stored in B+ tree? & Exclusively in leaf nodes \\
22 & Maximum B+ tree leaf capacity at height $h$, order $m$? & $(m-1)m^h$ \\
23 & Maximum edges in simple undirected graph with $V$ vertices? & $\frac{V(V-1)}{2}$ \\
24 & Sum of degrees in an undirected graph? & $2E$ (Handshaking Lemma) \\
25 & Maximum edges in complete bipartite graph $K_{a,b}$? & $a \times b$ \\
26 & Graph traversal algorithm implemented using Queue? & BFS (Breadth-First Search) \\
27 & Graph traversal algorithm implemented using Stack/Recursion? & DFS (Depth-First Search) \\
28 & Expected search complexity in a well-distributed hash table? & $O(1)$ \\
29 & Worst-case search complexity in a hash table? & $O(n)$ (all keys collide) \\
30 & Mathematical definition of hash table load factor $\alpha$? & $\alpha = n / m$ \\
31 & Collision resolution method prone to primary clustering? & Linear probing \\
32 & Purpose of tombstones during open addressing deletion? & Preserving unbroken probe search chains \\
33 & Space complexity of graph adjacency matrix? & $O(V^2)$ \\
34 & Space complexity of graph adjacency list? & $O(V + E)$ \\
35 & Graph representation providing instantaneous $O(1)$ edge test? & Adjacency matrix \\
36 & Representation most memory-efficient for sparse graphs? & Adjacency list \\
\end{longtable}
